%              --- User Research --- User Research --- User Research --- User Research --- User Research --- User Research --- User Research ---
\chapter{User Research - example middle chapter} % --- User Research --- User Research --- User Research --- User Research --- User Research ---
%              --- User Research --- User Research --- User Research --- User Research --- User Research --- User Research --- User Research ---

%       --- Methods --- Methods --- Methods --- Methods --- Methods --- Methods --- Methods --- Methods --- Methods --- Methods --- Methods --- Methods 
\section{Methods} % --- Methods --- Methods --- Methods --- Methods --- Methods --- Methods --- Methods --- Methods --- Methods --- Methods --- Methods 
%       --- Methods --- Methods --- Methods --- Methods --- Methods --- Methods --- Methods --- Methods --- Methods --- Methods --- Methods --- Methods 

\begin{tcolorbox}[      % Delete me later
    boxrule=0pt,        % Delete me later
    colback=white,      % Delete me later
    coltext=red         % Delete me later
]                       % Delete me later
    Describe what methods you have used to identify user needs. This can include methods such as storyboards and interviews and surveys with target users. Describe how you analyze the data you collect. Include an example. 
    
    \vspace{15pt}       % Delete me later
    \begin{center}      % Delete me later
    \textbf{Don't forget to delete this explanatory red text/block introduction, and replace it with your text in black.}
    \end{center}        % Delete me later
\end{tcolorbox}         % Delete me later

%       --- Stakeholder needs --- Stakeholder needs --- Stakeholder needs --- Stakeholder needs --- Stakeholder needs --- Stakeholder needs ---
\section{Stakeholder needs} % --- Stakeholder needs --- Stakeholder needs --- Stakeholder needs --- Stakeholder needs --- Stakeholder needs ---
%       --- Stakeholder needs --- Stakeholder needs --- Stakeholder needs --- Stakeholder needs --- Stakeholder needs --- Stakeholder needs ---

\begin{tcolorbox}[      % Delete me later
    boxrule=0pt,        % Delete me later
    colback=white,      % Delete me later
    coltext=red         % Delete me later
]                       % Delete me later
    Who are the stakeholders for this system? Provide a short description of what you know (so far) about each of the stakeholders and their needs. Highlight cases where their needs may differ. Personas could be appropriate here.
    
    Describe whose needs your system will prioritize. It is ok to state that a potential stakeholder is out of scope for the project. For example, a virtual tour of the SCU campus might target prospective students, but choose \textit{not} optimize for their parents (who are also potential stakeholders/viewers, but have different needs).

    \vspace{15pt}       % Delete me later
    \begin{center}      % Delete me later
    \textbf{Don't forget to delete this explanatory red text/block introduction, and replace it with your text in black.}
    \end{center}        % Delete me later
\end{tcolorbox}         % Delete me later

%       --- User stories --- User stories --- User stories --- User stories --- User stories --- User stories --- User stories --- User stories ---
\section{User stories} % --- User stories --- User stories --- User stories --- User stories --- User stories --- User stories --- User stories ---
%       --- User stories --- User stories --- User stories --- User stories --- User stories --- User stories --- User stories --- User stories ---

\begin{tcolorbox}[      % Delete me later
    boxrule=0pt,        % Delete me later
    colback=white,      % Delete me later
    coltext=red         % Delete me later
]                       % Delete me later
    For each stakeholder that you choose to prioritize, describe one or more user stories that your system will support, e.g., "As a prospective student, I want to find a social group at SCU so that I have friends to hang out with."
    
    Based on our user research and the priority stakeholders, we formulate the following user stories to guide system design. Each story captures a specific goal from a stakeholder's perspective and informs both feature planning and testing.
    
    \subsubsection*{Prospective Students}
    \begin{itemize}
      \item \textbf{As a prospective student,} I want to take a virtual tour of SCU’s dorms, \textbf{so that} I can compare housing options before visiting campus.
      \item \textbf{As a prospective student,} I want to preview clubs and student groups, \textbf{so that} I can plan how to get involved socially.
    \end{itemize}
    
    \subsubsection*{Parents of Prospective Students}
    \begin{itemize}
      \item \textbf{As a parent,} I want to access safety and health information during the tour, \textbf{so that} I feel confident in my child’s well-being on campus.
      \item \textbf{As a parent,} I want to know about academic support services, \textbf{so that} I can trust the school’s ability to help my child succeed.
    \end{itemize}
    
    \subsubsection*{Admissions Officers}
    \begin{itemize}
      \item \textbf{As an admissions officer,} I want to update featured content on the tour, \textbf{so that} we can highlight new programs each quarter.
      \item \textbf{As an admissions officer,} I want to track which parts of the tour are most visited, \textbf{so that} I can refine our outreach strategy.
    \end{itemize}
    
    \subsubsection*{Student Tour Guides}
    \begin{itemize}
      \item \textbf{As a student ambassador,} I want to record guided walkthroughs of campus spots, \textbf{so that} I can offer a personal touch remotely.
    \end{itemize}
    \vspace{15pt}       % Delete me later
    \begin{center}      % Delete me later
    \textbf{Don't forget to delete this explanatory red text/block introduction, and replace it with your text in black.}
    \end{center}        % Delete me later
\end{tcolorbox}         % Delete me later